\chapter{Conception globale de la solution}

\section{Introduction}
% TODO: Introduire la vision globale avant les chapitres techniques detailles.

\section{Vision generale de la solution}
% TODO: Presenter la solution comme une architecture data-driven en quatre couches.

\section{Architecture globale}
\begin{figure}[htbp]
    \centering
    \fbox{\parbox{0.82\textwidth}{\centering Espace reserve au schema d'architecture globale de la solution.}}
    \caption{Architecture globale de la solution data-driven}
    \label{fig:architecture_globale}
\end{figure}

\section{Separation des couches fonctionnelles}

\subsection{Couche Data Engineering}
% TODO: Presenter le role de l'ETL, de MySQL, de fact_coverage et de survey_responses.

\subsection{Couche Business Intelligence}
% TODO: Presenter le role des dashboards Power BI pour le management et les equipes internes.

\subsection{Couche application mobile client}
% TODO: Presenter l'application mobile comme interface client Unilever orientee magasins affectes.

\subsection{Couche back-office validation interne}
% TODO: Presenter le portail interne Smollan pour les anomalies et la validation.

\section{Choix technologiques}
% TODO: Justifier les choix Python, MySQL, Power BI, Spring Boot, React et Docker.

\section{Conclusion}
% TODO: Faire la transition vers le pipeline Data Engineering.

